\section{本讲问题：把 scaling laws 当成工程预测工具}


\begin{knowledgebox}{讲义补充：源 Slide 1 标记课程转向}
前面几讲已经讨论数据、GPU、kernel 和并行训练。Lecture 9 关注一个更上层的问题：当你真的有大规模算力时，如何决定模型大小、数据量、训练步数、架构、优化器和超参数。Scaling laws 的角色是把这些决定从经验崇拜变成可预测的工程决策。
\end{knowledgebox}


\singleslide{slides-images/slide-002.jpg}{“一万张 B200 用一个月”的资源决策：基础设施和数据准备完成后，究竟训练多大的模型。}{2}

Slide 2 把 scaling law 从拟合曲线变成资本配置问题。硬件时间固定、训练失败代价巨大，团队必须在正式 run 前用小模型预测 architecture、parameter count、token budget 与训练时长；否则再好的 infra 和 dataset 也可能被错误的模型规模浪费。

\begin{knowledgebox}{讲义补充：源 Slide 2 的场景为什么尖锐}
这页把 scaling laws 放到真实资源决策里。算力巨大时，随机试错非常昂贵；照抄已有模型配置也不可靠，因为硬件、数据、目标和预算不同。你需要先设计小规模实验，再预测大规模训练结果，最后决定把训练 compute 放在哪里。
\end{knowledgebox}

\begin{figure}[H]
\centering
\includegraphics[width=0.86\textwidth]{slide-003.jpg}
\caption{Slide 3：Scaling 并不容易，架构宽深、heads、nonlinearity 等选择不能只靠 cargo cult。}
\figsource{CS336 Spring 2026 Lecture 9 官方幻灯片。}
\end{figure}

\begin{warningbox}{读图：Slide 3 反对的是盲目照抄}
公开模型给了很多配置，但它们是特定资源、数据、训练目标和组织能力下的结果。Scaling law 实验要回答的是：哪些选择能在小规模上拟合出稳定趋势，并且外推到目标规模时仍可信。它不是魔法，也不是替代所有大规模验证。
\end{warningbox}

\begin{figure}[H]
\centering
\includegraphics[width=0.86\textwidth]{slide-004.jpg}
\caption{Slide 4：本讲目标是用简单、可预测的 laws 描述 LM 行为，从大模型调参转向小模型外推。}
\figsource{CS336 Spring 2026 Lecture 9 官方幻灯片。}
\end{figure}

\begin{importantbox}{术语消化：什么是 scaling law}
Scaling law 是把某个规模变量映射到 loss、error 或 performance 的经验公式。常见形式是 power law，例如
\[
L(x)=A x^{-\alpha}+C,
\]
其中 \(x\) 可以是数据量、模型参数量或计算量，\(A\) 是尺度系数，\(\alpha\) 是 scaling exponent，\(C\) 是不可约误差或 offset。它的工程价值在于：用小规模点估计曲线，再预测大规模资源分配。
\end{importantbox}

